\documentclass[11pt]{article}

\usepackage[T1]{fontenc}
\usepackage{amsmath}
\usepackage{amssymb}
\usepackage{booktabs}
\usepackage[margin=1in]{geometry}
\usepackage{xcolor}
\usepackage{listings}
\usepackage[hidelinks]{hyperref}

\title{VBOGS Algorithm}
\author{VBOGS}
\date{}

\newcommand{\R}{\mathbb{R}}
\newcommand{\E}{\mathbb{E}}
\newcommand{\Umax}{U_{\max}}
\DeclareMathOperator{\round}{round}

\definecolor{codebg}{RGB}{248,248,248}
\definecolor{codeframe}{RGB}{210,210,210}
\definecolor{codekw}{RGB}{35,80,150}
\definecolor{codecomment}{RGB}{90,105,90}
\definecolor{codenum}{RGB}{120,120,120}

\lstdefinelanguage{VBOGSPseudo}{
  sensitive=true,
  morekeywords={
    Function,Return,Foreach,For,If,Else,While,Break,continue,in,range,None,
    zeros,concat,copy,mean,chunks,defaultdict,zip,min
  },
  morecomment=[l]{\#}
}

\lstset{
  language=VBOGSPseudo,
  basicstyle=\ttfamily\small,
  keywordstyle=\bfseries\color{codekw},
  commentstyle=\itshape\color{codecomment},
  numbers=left,
  numberstyle=\scriptsize\color{codenum},
  numbersep=8pt,
  frame=single,
  framerule=0.25pt,
  rulecolor=\color{codeframe},
  backgroundcolor=\color{codebg},
  breaklines=true,
  columns=fullflexible,
  keepspaces=true,
  showstringspaces=false,
  tabsize=4,
  xleftmargin=1.5em,
  framexleftmargin=1.2em
}

\begin{document}
\maketitle

\section{Purpose}

VBOGS combines Octree-AnyGS with a per-anchor Variational Bayes Gaussian
Splatting (VBGS) uncertainty head. Octree-AnyGS provides the primary scene
representation. VBGS provides calibrated posterior uncertainty over the points
assigned to each Octree-AnyGS anchor. The final output is a next-best-view
camera pose chosen to maximize expected posterior uncertainty over visible
surface.

\begin{table}[h]
  \centering
  \footnotesize
  \begin{tabular}{@{}p{0.44\linewidth}p{0.46\linewidth}@{}}
    \toprule
    Component & Role \\
    \midrule
    Primary scene representation &
      Octree-AnyGS trained on posed RGB frames in PyTorch \\
    Stereo depth cloud &
      Used to produce 6D point observations for the uncertainty head \\
    Uncertainty head &
      Per-anchor Variational Bayes GMM from VBGS in JAX \\
    Decision output &
      Next-best camera pose for uncertainty reduction \\
    \bottomrule
  \end{tabular}
\end{table}

\section{Terminology}

Octree-AnyGS stores anchors as a flat tensor. Each anchor has a level
\(\ell\), and its cell size is

\[
  cell\_size(\ell)
  =
  \frac{\mathrm{voxel\_size}}{\mathrm{fork}^{\ell}}.
\]

There is no separate leaf voxel object. In this algorithm, an anchor is the
voxel at its level of detail, so ``anchor'' and ``voxel'' are used
interchangeably.

\section{Hyperparameters}

\begin{table}[h]
  \centering
  \begin{tabular}{@{}p{0.30\linewidth}p{0.62\linewidth}@{}}
    \toprule
    Hyperparameter & Meaning \\
    \midrule
    \texttt{K\_INIT} &
      Initial number of GMM components per anchor; default is \texttt{n\_offsets}. \\
    \texttt{K\_MAX} &
      Maximum number of components per anchor. \\
    \texttt{K\_GROWTH\_FACTOR} &
      Rule for increasing \(K\) when an anchor is under-modeled, e.g. \(K \mapsto 2K\). \\
    \texttt{ELBO\_IMPROVEMENT\_TOL} &
      Minimum per-point ELBO gain required to accept a larger \(K\). \\
    \texttt{VBEM\_BATCH\_SIZE} &
      Number of points per VBEM sufficient-statistics batch. \\
    \texttt{MIN\_POINTS\_PER\_ANCHOR} &
      Anchors below this point count are flagged as unobserved and receive \(\Umax\). \\
    \texttt{NBV\_CANDIDATES} &
      Candidate camera poses to score. \\
    \bottomrule
  \end{tabular}
\end{table}

\section{Stage 1: Train the Scene}

Stage 1 is unchanged from upstream Octree-AnyGS. The scene is trained from
posed RGB stereo frames. Stereo is not used as a training signal for this
stage; it is used later to build the 6D point cloud.

\begin{lstlisting}
scene = Octree_AnyGS.train(input_data, gaussian_type=3DGS)

# input_data =
# {
#   (cam_pose_L, cam_pose_R, img_L, img_R, K_L, K_R, baseline)
# }
\end{lstlisting}

\section{Stage 2: Build a 6D Stereo Point Cloud}

For each rectified stereo pair, estimate disparity, convert disparity to depth,
unproject into the left-camera frame, transform to world coordinates, attach
RGB color, and concatenate all frames into one 6D point cloud:

\[
  x_n^{w}
  =
  \begin{bmatrix}
    p_n^{w} \\
    c_n
  \end{bmatrix}
  \in \R^6,
  \qquad
  p_n^{w} \in \R^3,
  \qquad
  c_n \in \R^3.
\]

\begin{lstlisting}
points_world = []                                # (N, 6) = xyz + rgb
Foreach (cam_L, cam_R, img_L, img_R, K_L, K_R, baseline) in input_data:
    disparity = StereoMatch(img_L, img_R)        # SGBM / RAFT-Stereo / etc.
    depth_L   = (K_L.fx * baseline) / disparity  # mask invalid disparities
    pts_cam   = Unproject(depth_L, K_L)          # (H*W, 3), left-cam frame
    pts_world = cam_L.transform(pts_cam)         # world frame
    rgb       = img_L.reshape(-1, 3)
    points_world.append(concat(pts_world, rgb, axis=-1))

points_world = concat(points_world)              # (sum_i N_i, 6)

# VBGS expects normalized data.
points_norm, norm_params = normalize_data(points_world)
\end{lstlisting}

Both coordinate frames are retained:

\begin{itemize}
  \item \texttt{points\_world} is used for bucketing, because Octree-AnyGS
        anchors live in world coordinates.
  \item \texttt{points\_norm} is used for VBGS fitting, because the VBGS priors
        expect normalized inputs.
\end{itemize}

\section{Stage 3: Per-Anchor VBGS Fitting}

\subsection{Anchor Grid}

Point-to-anchor assignment uses the same grid discretization that Octree-AnyGS
uses when it creates anchors. A point and anchor match at level \(\ell\) when
their integer grid coordinates match:

\[
  cur\_size(\ell)
  =
  \frac{\texttt{scene.voxel\_size}}{\texttt{scene.fork}^{\ell}},
\]

\[
  grid\_coord(p,\ell)
  =
  \round
  \left(
    \frac{p - \texttt{scene.init\_pos}}{cur\_size(\ell)}
  \right).
\]

For an anchor \(a\),

\[
  grid\_coord(a,\ell_a)
  =
  \round
  \left(
    \frac{a_{\mathrm{pos}} - \texttt{scene.init\_pos}}
         {cur\_size(\ell_a)}
  \right).
\]

The padding offset that Octree-AnyGS adds when materializing anchor positions
cancels out because both points and anchors are converted back to grid
coordinates. Do not subtract padding during bucketing.

\subsection{Real Scene Fields}

The implementation reads fields directly from the Octree-AnyGS scene object:

\begin{lstlisting}
# pc = scene
# Real fields:
#   pc._anchor      [N, 3], world coordinates
#   pc._level       [N, 1], integer level
#   pc.voxel_size
#   pc.fork
#   pc.init_pos
#   pc.n_offsets
# There is no get_anchors() method.

anchor_pos   = pc._anchor
anchor_level = pc._level.squeeze(1)
\end{lstlisting}

\subsection{Anchor Lookup}

Precompute a lookup from level and grid coordinate to anchor id:

\begin{lstlisting}
anchor_index = {}   # (level, grid_coord_tuple) -> anchor_id
For i in range(N):
    l  = anchor_level[i]
    gc = round((anchor_pos[i] - pc.init_pos) / cur_size(l))
    anchor_index[(l, tuple(gc))] = i
\end{lstlisting}

\subsection{Assign Points to Every Containing Level}

Each point is assigned to every anchor that contains it across all Octree-AnyGS
levels, not only to the finest matching anchor. This is important because
Octree-AnyGS chooses levels of detail at render time. A coarse anchor must
therefore receive the points it would represent if it were selected by a far
camera view.

\begin{lstlisting}
pts_by_anchor = defaultdict(list)

For p_world, p_norm in zip(points_world, points_norm):
    For l in range(scene.levels):                 # all levels, not just finest
        gc = tuple(round((p_world[:3] - pc.init_pos) / cur_size(l)))
        If (l, gc) in anchor_index:
            pts_by_anchor[anchor_index[(l, gc)]].append(p_norm)
\end{lstlisting}

\subsection{Fit One Anchor}

For each observed anchor, fit a VBGS mixture over that anchor's normalized 6D
points. The fit proceeds in VBEM batches using the original VBGS update
functions.

\begin{lstlisting}
Function FitAnchor(pts_a, K):
    mean_init   = init_means(pts_a, K)           # k-means++ or random_mean_init
    prior_model = get_volume_delta_mixture(
        n_components=K,
        mean_init=mean_init,
        position_event_shape=(3, 1),
        beta=0,
        learning_rate=1)

    model = copy(prior_model)
    prior_stats = space_stats = color_stats = None

    For each batch in chunks(pts_a, VBEM_BATCH_SIZE):
        model, prior_stats, space_stats, color_stats = fit_gmm_step(
            prior_model,
            model,
            batch,
            VBEM_BATCH_SIZE,
            prior_stats,
            space_stats,
            color_stats)

    elbo_arr, _    = compute_elbo_delta(model, pts_a)   # (N,)
    elbo_per_point = mean(elbo_arr)
    Return model, elbo_per_point
\end{lstlisting}

\subsection{Grow the Component Count}

The mixture starts at \(\texttt{K\_INIT}\). Larger component counts are accepted
only when they improve mean per-point ELBO enough:

\[
  \Delta \bar{\mathcal{L}}
  =
  \bar{\mathcal{L}}_{K_{\mathrm{next}}}
  -
  \bar{\mathcal{L}}_K
  \geq
  \texttt{ELBO\_IMPROVEMENT\_TOL}.
\]

\begin{lstlisting}
anchor_posterior = [None] * N

For i in range(N):
    pts_a = pts_by_anchor[i]

    If |pts_a| < MIN_POINTS_PER_ANCHOR:
        anchor_posterior[i] = None               # unobserved -> high uncertainty
        continue

    K = K_INIT                                   # default: scene.n_offsets
    model, elbo = FitAnchor(pts_a, K)

    While K < K_MAX:
        K_next              = min(K_MAX, K_GROWTH_FACTOR(K))
        model_next, elbo_nx = FitAnchor(pts_a, K_next)

        If elbo_nx - elbo < ELBO_IMPROVEMENT_TOL:
            Break                                # adding components no longer helps

        model, elbo, K = model_next, elbo_nx, K_next

    anchor_posterior[i] = model

    # Optional: if final K hit K_MAX with ELBO still improving, flag anchor
    # as under-modeled and let Octree-AnyGS densify it on a follow-up pass.
\end{lstlisting}

The ELBO criterion is pragmatic rather than fully principled: the KL term
scales with component count, so mean per-point ELBO is biased across \(K\).
Held-out log likelihood or BIC would be cleaner model-selection alternatives.

\section{Stage 4: Reduce Posterior to a Scalar Uncertainty}

Each fitted component has a Normal-Wishart posterior over its spatial Gaussian
parameters and a delta-color posterior for RGB. The mixture prior is
Dirichlet-distributed with parameters \(\alpha\). The uncertainty scalar for
an observed anchor is the expected posterior entropy, weighted by expected
mixture weights:

\[
  \bar{\pi}_{ak}
  =
  \E_q[\pi_{ak}]
  =
  \frac{\alpha_{ak}}{\sum_j \alpha_{aj}},
\]

\[
  U_a
  =
  \sum_{k=1}^{K_a}
  \bar{\pi}_{ak}
  \left(
    H_{\mathrm{NW}}(q^s_{ak})
    +
    H_{\Delta}(q^c_{ak})
  \right).
\]

\begin{lstlisting}
U = zeros(N)

For i in range(N):
    If anchor_posterior[i] is None:
        U[i] = U_MAX                             # unseen anchor = maximally uncertain
    Else:
        q      = anchor_posterior[i]
        pi_bar = E[q.mixture.prior]              # (K,) expected mixing weights
        H_k    = NormalWishartEntropy(q.likelihood) + DeltaEntropy(q.delta)
        U[i]   = sum_k(pi_bar[k] * H_k[k])       # expected posterior entropy

        # Entropies are computed in normalized coordinates, so comparisons
        # across anchors are unit-consistent.
\end{lstlisting}

\section{Stage 5: Next-Best-View Selection}

Stage 5 projects per-anchor uncertainty through Octree-AnyGS's own
LOD-selected splatting path. This produces a pixel-space uncertainty image for
each candidate view and automatically respects visibility, occlusion, and LOD.

Octree-AnyGS normally obtains colors from \texttt{get\_color\_mlp()} and does
not expose an auxiliary anchor channel. VBOGS reuses the geometry path from
\texttt{generate\_gaussians}, substitutes a scalar uncertainty color, and then
rasterizes.

\subsection{Scalar Rendering}

\begin{lstlisting}
Function render_scalar(cam, pc, per_anchor_scalar):
    # 1. LOD mask selection, same as render().
    #    At NBV time Octree-AnyGS is past training. If pc.progressive is False,
    #    the iteration argument is ignored. Passing a large iteration, or using
    #    set_anchor_mask_perlevel with cur_level = pc.levels - 1, forces all
    #    levels to be active.
    pc.set_anchor_mask(
        cam.camera_center,
        iteration=INT_MAX,
        resolution_scale=cam.resolution_scale)

    visible_mask = pc._anchor_mask

    # 2. Reuse geometry, opacity, and MLP selection mask from the normal path.
    xyz, _color, opacity, scaling, rot, _sh, selection_mask = \
        pc.generate_gaussians(cam, visible_mask)

    # 3. Broadcast anchor uncertainty over offsets, then apply the same
    #    opacity-MLP selection mask used by generate_gaussians.
    u_per_anchor   = per_anchor_scalar[visible_mask]                # [M]
    u_per_offset   = u_per_anchor.repeat_interleave(pc.n_offsets)   # [M*n_offsets]
    u_per_gaussian = u_per_offset[selection_mask][:, None]          # [G, 1]

    # 4. Rasterize with zero background so empty pixels contribute zero to both
    #    uncertainty and alpha images.
    unc_image, alpha_image, _ = gsplat.rasterization(
        means=xyz,
        quats=rot,
        scales=scaling,
        opacities=opacity.squeeze(-1),
        colors=u_per_gaussian,
        viewmats=cam.world_view_transform.T[None],
        Ks=cam.K[None],
        backgrounds=zeros(1),
        width=cam.W,
        height=cam.H,
        packed=False,
        sh_degree=None,                          # raw scalar, not SH
        render_mode=pc.render_mode)

    Return unc_image[0, ..., 0], alpha_image[0, ..., 0]   # both (H, W)
\end{lstlisting}

\subsection{Candidate Scoring}

Candidate views are scored by uncertainty per unit of visible surface:

\[
  score(cam)
  =
  \frac{\sum_{u,v} I_U(u,v)}
       {\sum_{u,v} I_\alpha(u,v) + \epsilon}.
\]

This alpha-normalized score rewards poses that see uncertain content rather
than poses that simply see more content. An optional travel-cost penalty can
be subtracted when reachability or motion budget matters.

\begin{lstlisting}
best_score, best_pose = -inf, None

Foreach cam in NBV_CANDIDATES:
    unc_image, alpha_image = render_scalar(cam, scene, U)

    score = sum(unc_image) / (sum(alpha_image) + EPS)

    # Optional motion or reachability term:
    # score -= lambda * TravelCost(current_pose, cam)

    If score > best_score:
        best_score, best_pose = score, cam

Return best_pose
\end{lstlisting}

\section{Implementation Notes}

\begin{itemize}
  \item \textbf{Framework bridge.} Octree-AnyGS runs in PyTorch and VBGS runs
        in JAX. The framework boundary is the filesystem: pass NumPy arrays
        across processes rather than sharing tensors or autograd graphs.
  \item \textbf{Color model.} VBGS uses raw RGB with a delta MVN and fixed
        precision. Octree-AnyGS uses spherical harmonics. This is acceptable
        because VBOGS uses VBGS posterior parameters for uncertainty, not VBGS
        renders.
  \item \textbf{Stereo quality.} Stereo depth quality dominates
        anchor-point-count sparsity at range. Mask low-confidence disparities
        before unprojection using checks such as left-right consistency and
        texture thresholds.
  \item \textbf{Continual variant.} The streaming update path uses
        \path{fit_gmm_step}. It can carry \path{prior_stats},
        \path{space_stats}, and \path{color_stats} across batches when
        updating the uncertainty head online.
  \item \textbf{Empty-region blindness.} \texttt{render\_scalar} only splats
        through existing anchors. Truly unobserved regions with no anchor at
        any LOD contribute zero to the NBV score, so this formulation directs
        the camera toward poorly modeled observed content rather than never-seen
        space. Exploration of empty space would require an occupancy prior or
        unknown-ray penalty as a follow-on extension.
\end{itemize}

\end{document}
